% !TEX encoding = UTF-8
% !TEX TS-program = pdflatex
% !TEX root = ../tesi.tex

%**************************************************************
\chapter{Descrizione dello stage}
\label{cap:processi-metodologie}
%**************************************************************

\intro{In questo capitolo verrà descritto in dettaglio lo stage, l'azienda dove è stato svolto e come è stato organizzato il lavoro.}

%**************************************************************

\section{L'azienda}

Sync Lab S.r.l. è un'azienda di consulenza informatica nata nel 2002 nella sede di Napoli e trasformatasi molto velocemente nel corso degli anni in System Integrator grazie ad un processo di maturazione delle competenze tecnologiche. Supporta le esigenze di innovazione dei clienti offrendo soluzioni IT in ambito Business Consultancy, Project Financing e IT Consultancy.
\begin{figure}[h]
	\begin{center}
		\includegraphics[width=7cm]{logo_large}
		\caption{Logo di Sync Lab}
	\end{center}
\end{figure}
Ad oggi l'azienda conta un organico di circa 200 dipendenti ed è riuscita a coprire tutto il territorio nazionale fino ad ottenere un totale di cinque sedi a Roma, Napoli, Verona, Padova e Milano.\\
L'azienda, propone sul mercato prodotti software, attraverso essi Sync Lab ha gradualmente conquistato significativamente fette di mercato nei seguenti settori: mobile, videosorveglianza e sicurezza delle infrastrutture informatiche aziendali.
Il loro obiettivo é quello di supportare il cliente nella Realizzazione, Messa in Opera e Governance di soluzioni IT, sia dal punto di vista Tecnologico, sia nel Governo del Cambiamento Organizzativo.

%**************************************************************
\section{L'idea dello stage}

Il progetto da svolgere proposto dall'azienda per lo stage è una web application chiamata NFTLab legata all'ambito blockchain e NFT, due argomenti che stanno spopolando al momento e quindi è uno stage figo perchè tratta di cose contemporanee.\\
In dettaglio quando l'utente accede al sito potrà visualizzare le opere multimediali in vendita, fare l'accesso o la registrazione al sito. Dopo essere acceduto, l'utente potrà caricare le proprie opere e decidere se metterle in vendita o meno. Inoltre potrà sfogliare il catalogo delle opere presenti e decidere di comprarne a sua volta.\\
Essendo legato al concetto di blockchain e NFT le azioni di compravendita verranno effettuate tramite lo scambio di una moneta virtuale detta Ethereum. Un utente nel momento in cui carica un'opera diventa l'autore ed il proprietario di quest'ultima, ma se la venderà l'acquirente diventerà il nuovo proprietario ma solo l'utente che ha caricato l'opera è e sarà sempre l'autore.\\
Per realizzare questi aspetti sono state implementate diverse maschere, ovvero interfacce utente, tramite il framework Vue.js, uno strumento utile allo sviluppo di web appication.\\
Parte del percorso di stage è stata riservata ad un periodo di studio (senza implementazioni se non a scopo didattico) delle tecnologie legate al back end , in particolare il framework Spring.

%**************************************************************
\section{Organizzazione del lavoro}

L'azienda adotta il modello di sviluppo Agile che si basa sull'interazione continua con gli stakeholder e predilige:
\begin{itemize}
	\item gli individui e le loro interazioni piuttosto che i processi e gli strumenti;
	\item il software funzionante piuttosto che una documentazione esaustiva;
	\item la collaborazione con il cliente piuttosto che la negoziazione dei contratti;
	\item essere propositivi verso il cambiamento piuttosto che rifiutarlo.
\end{itemize}
L'idea di questo modello non si basa quindi su uno sviluppo sequenziale ma sul concetto di rivedere di continuo le specifiche adeguandole durante l'avanzamento dello sviluppo software. In questo modo si possono apportare agilmente modifiche mediante metodologie iterative ed incrementali ed un continuo scambio di pareri ed informazioni tra gli sviluppatori ed il committente. Il progetto di stage è stato portato avanti utilizzando il metodo Scrum.\\
Scrum è il metodo più diffuso e prevede di dividere il progetto in blocchi di lavoro.

%**************************************************************